
@@M0@@ ; on a le diagramme commutatif :
@@M1@@
où les morphismes @@M2@@ et @@M3@@ sont surjectifs. Soient @@M4@@, pour @@M5@@, des éléments de @@M6@@ tels que pour tout @@M7@@, on ait
@@M8@@
On définit un homomorphisme @@M9@@ en posant @@M10@@. On obtient ainsi une immersion fermée @@M11@@ qui est une section de @@M12@@, et qui est telle que @@M13@@. De plus, étant une section d'un morphisme lisse, c'est une immersion régulière (VII 1.10). Sachant que @@M14@@ et @@M15@@ sont régulières, on en déduira que @@M16@@ est régulière si on montre que la suite :
@@M17@@
est exacte et localement scindée, @@M18@@, @@M19@@ et @@M20@@ désignant respectivement les idéaux définissant @@M21@@, @@M22@@, et @@M23@@ (VII 1.7). Il suffit pour cela de montrer que @@M24@@ est facteur direct de @@M25@@. Or on voit tout de suite que @@M26@@ se déduit de @@M27@@ par le changement de base @@M28@@. En outre, les immersions @@M29@@ et @@M30@@ sont régulières (VII 1.10). Par conséquent (VII 1.2), l'idéal qui définit @@M31@@ est @@M32@@, et le faisceau conormal de @@M33@@ est @@M34@@. La suite exacte des faisceaux conormaux relative à la décomposition @@M35@@ donne donc en particulier un homomorphisme : @@M36@@. Il est immédiat de vérifier que le composé : @@M37@@ est l'identité (du point de vue affine, cela correspond à la décomposition de @@M38@@ en somme directe d'un idéal de @@M39@@, et de l'idéal engendré par les @@M40@@). On a donc bien ainsi le scindage cherché.

